% =====================================================================
%  Response to Reviewer 5
%  Manuscript: Multi-sensor dataset of a tendon-driven continuum robot
%               in dynamic motion
%
%  This file compiles standalone (pdflatex twice).
%  To merge it into a larger response letter, delete everything from
%  \documentclass down to \begin{document}, delete the closing
%  \end{document} and the thebibliography block, and \input this file.
% =====================================================================
\documentclass[11pt,a4paper]{article}
\usepackage[a4paper,margin=1in]{geometry}
\usepackage{amsmath}
\usepackage{amssymb}
\usepackage{siunitx}
\usepackage{xcolor}
\usepackage{parskip}
\usepackage{hyperref}

\bibliographystyle{ieeetr}

% --- helper macros -----------------------------------------------------
% Reviewer text is typeset in italics so it is visually separated from
% the authors' response.
\newcommand{\revcomment}[1]{\textit{#1}}
% Placeholder for sections still to be written. Delete as you fill them in.
\newcommand{\todoresp}{\textcolor{red}{[To be completed.]}}
% -----------------------------------------------------------------------

\begin{document}

\section{reviewer 5}

\paragraph{general comment}
\revcomment{The work presents a dataset for a single-segment soft robot, built from commercially
available components. The dataset contains readings from six sensor modalities in quasi-static
and dynamic motions as well as contact situation. Cross-validation was undertaken for
consistency assessment. Overall, the work addresses a general lack of comparable,
transferable, standardized and available experimental data in the field of soft robotics. The
scientific contribution and relevance of the work is well motivated and aims at filling an
existing research gap in the field of soft robotics. However, the paper and the dataset show
major weaknesses and inconsistencies, which should be revised before the paper can be
published. Details below.}

\paragraph{Response to the general comment}
\todoresp

% =====================================================================
\subsection{Major comments (paper)}
% =====================================================================

\paragraph{Comment 5.1}
\revcomment{The introduction to the work is thoroughly written and also effectively compares the
state of the art in soft robotics well with other data-intensive research fields such as
autonomous driving and robotics in general. However, the introduction lacks comparisons to
datasets in the same field, which should be added/referred to, e.g., \cite{Hansen2026,
Muhmann2025, Zhang2025, Liu2026, Wang2025}.}

\paragraph{Response to Comment 5.1}
\todoresp

\paragraph{Changes made to the article}
\todoresp

\paragraph{Comment 5.2}
\revcomment{While the experimental setup is described in detail, the figures (especially Fig.~2)
should be much clearer:
\begin{itemize}
  \item Inconsistent writing (lowercase / uppercase). Example: Contact probe vs.\ Reflective Marker.
  \item Coordinate frames are barely visible, especially when printed on paper in greyscale.
        Increase contrast.
  \item Dashed lines and circles are not aligned well. Correct alignment.
\end{itemize}}

\paragraph{Response to Comment 5.2}
\todoresp

\paragraph{Changes made to the article}
\todoresp

\paragraph{Comment 5.3}
\revcomment{Multiple qualitative terms such as ``rigidly connected'' and ``geometrically stable''
are used. Such vague terms are too unprecise in an experimental paper about soft robotics,
which heavily relies on a sufficiently rigid mounting structure. The authors should provide
quantitative measurements such as static directional stiffness values or structural
eigenfrequencies to quantify that the mounting structure is sufficiently stiff for the
experiments.}

\paragraph{Response to Comment 5.3}
\todoresp

\paragraph{Changes made to the article}
\todoresp

\paragraph{Comment 5.4}
\revcomment{The contact-force validation remains largely qualitative. Terms such as ``close
agreement'' and ``good agreement'' should be supported by quantitative metrics such as RMSE,
mean bias, maximum absolute error, and preferably correlation values.}

\paragraph{Response to Comment 5.4}
\todoresp

\paragraph{Changes made to the article}
\todoresp

\paragraph{Comment 5.5}
\revcomment{Inconsistency in Tab.~1: The table does not only contain geometrical parameters.
Correct the description of the table.}

\paragraph{Response to Comment 5.5}
\todoresp

\paragraph{Changes made to the article}
\todoresp

\paragraph{Comment 5.6}
\revcomment{Inconsistency in Tab.~1: The physical unit of weight is newtons (N), not grams (g).
Refer to mass, not weight, when using grams as the physical unit. Correct units.}

\paragraph{Response to Comment 5.6}
\todoresp

\paragraph{Changes made to the article}
\todoresp

\paragraph{Comment 5.7}
\revcomment{Inconsistency with kinematic definitions: $g_{\mathrm{fix}}$ is introduced as a
rotation matrix (orientation), but $g_{\mathrm{wand}}$ is introduced as a pose. In robotics, a
pose is position and orientation.}

\paragraph{Response to Comment 5.7}
\todoresp

\paragraph{Changes made to the article}
\todoresp

\paragraph{Comment 5.8}
\revcomment{The paper currently lacks a clear, dedicated conclusion and summary section. The
authors should add a concluding section that summarizes the impact of the dataset and outlines
potential future use cases.}

\paragraph{Response to Comment 5.8}
\todoresp

\paragraph{Changes made to the article}
\todoresp

\paragraph{Comment 5.9}
\revcomment{The paper currently lacks an application to a baseline model for at least one of the
use cases that are mentioned in the introduction. Therefore, the paper does not demonstrate by
itself that the dataset is useful for any of these applications. Neither does it demonstrate
that the authors have successfully used their dataset. This could have helped with the
inconsistencies of the dataset (see further below).}

\paragraph{Response to Comment 5.9}
\todoresp

\paragraph{Changes made to the article}
\todoresp

\paragraph{Comment 5.10}
\revcomment{The GVS model used for Figure~6 and Table~7 is not included in the released code.
Since this model forms part of the technical validation, the implementation and all parameters
required to reproduce the reported results should be released.}

\paragraph{Response to Comment 5.10}
\todoresp

\paragraph{Changes made to the article}
\todoresp

% =====================================================================
\subsection{Minor comments (paper)}
% =====================================================================

\paragraph{Comment 5.11}
\revcomment{The word ``minimize'' should only be used when referring specifically to the results
of an optimization routine. Consider using ``reduce'' or ``eliminate'' to avoid confusion. This
occurs multiple times in the paper and should be corrected.}

\paragraph{Response to Comment 5.11}
\todoresp

\paragraph{Changes made to the article}
\todoresp

\paragraph{Comment 5.12}
\revcomment{Sec.~2.2: 3D printed pulleys appear to be an unnecessarily low-cost option for
routing cables. Have the authors checked for roundness (part may have a seam from printing) and
eccentricity after pressing the bearing in? It remains unclear why 3D printed pulleys were used
at all, as machined parts are widely available standard components.}

\paragraph{Response to Comment 5.12}
\todoresp

\paragraph{Changes made to the article}
\todoresp

\paragraph{Comment 5.13}
\revcomment{While mathematically identical for orthogonal matrices, it is rather common in
robotics to use the transpose of a rotation matrix instead of its inverse (see inversion of
$g_{\mathrm{base}}$ on page 16). The authors may want to use the transpose for stylistic
consistency.}

\paragraph{Response to Comment 5.13}
\todoresp

\paragraph{Changes made to the article}
\todoresp

\paragraph{Comment 5.14}
\revcomment{Some parts of the paper read more like a technical documentation rather than a
scientific paper. Overall, the writing could be improved before publication so that the text
reads more fluently.}

\paragraph{Response to Comment 5.14}
\todoresp

\paragraph{Changes made to the article}
\todoresp

\paragraph{Comment 5.15}
\revcomment{Data availability and code availability refer to the same internet source, so it
could be included in a single section rather than two short sections.}

\paragraph{Response to Comment 5.15}
\todoresp

\paragraph{Changes made to the article}
\todoresp

\paragraph{Comment 5.16}
\revcomment{Table~3 states that the FBG shape contains 502 points but uses the index range
0--502, which would correspond to 503 points.}

\paragraph{Response to Comment 5.16}
\todoresp

\paragraph{Changes made to the article}
\todoresp

% =====================================================================
\subsection{Major comments (dataset)}
% =====================================================================

\paragraph{Comment 5.17}
\revcomment{The documented column order of \texttt{mocap\_frames.csv} appears to be incorrect.
Table~3 states $[x, y, z, \mathrm{roll}, \mathrm{pitch}, \mathrm{yaw}]$, whereas both the
released processing code (\texttt{data\_optitrack.m}, lines 143--149;
\texttt{process\_data.m}, lines 1049--1066) and inspection of the released CSV indicate
$[\mathrm{roll}, \mathrm{pitch}, \mathrm{yaw}, x, y, z]$. This can directly lead to incorrect
interpretation of the dataset and should be corrected before publication.}

\paragraph{Response to Comment 5.17}
\todoresp

\paragraph{Changes made to the article}
\todoresp

\paragraph{Comment 5.18}
\revcomment{A Python logic error is present in \texttt{read4MotorCircle.py}:
\texttt{if motion\_pattern == "circle" or "Lissajous":}. This condition is always true because
the non-empty string literal \texttt{"Lissajous"} evaluates to true in a Boolean context in
Python. As a consequence, the angular tolerance is apparently set to $12^{\circ}$ for all motion
types rather than the intended $0.5^{\circ}$ default. The authors should clarify whether this
code version was used for the published experiments and assess its impact on the recorded
trajectories.}

\paragraph{Response to Comment 5.18}
\todoresp

\paragraph{Changes made to the article}
\todoresp

\paragraph{Comment 5.19}
\revcomment{The filtering procedure may be seriously affected by irregular timestamps.
\texttt{process\_data.m} determines the sampling frequency using
\texttt{Fs = 1/median(diff(t))} and applies a standard Butterworth filter before uniform
resampling. This assumes approximately equidistant samples. The inspected ATI raw recording
exhibits strongly burst-like host timestamps despite an average data rate of approximately
\SI{200}{\hertz}, meaning that the median timestamp interval may yield a grossly incorrect
filter sampling frequency. The authors should verify whether the released processed signals were
affected and, if necessary, regenerate them using a regular time grid before filtering.}

\paragraph{Response to Comment 5.19}
\todoresp

\paragraph{Changes made to the article}
\todoresp

\paragraph{Comment 5.20}
\revcomment{Contact-data processing is inconsistent between the manuscript and the released code.
The manuscript specifies a \SI{15}{\hertz} low-pass filter, whereas
\texttt{compare\_ft\_sensors.m} uses \SI{30}{\hertz}. Processed files are described as
header-less, while the contact script uses \texttt{writetable}. File names also differ
(\texttt{tendon\_tensions.csv} vs.\ \texttt{cable\_tensions.csv},
\texttt{contact\_wrench.csv} vs.\ \texttt{wrench\_wand.csv}), and no clear generator for the
documented \texttt{wand\_pose.csv} was identified.}

\paragraph{Response to Comment 5.20}
\todoresp

\paragraph{Changes made to the article}
\todoresp

\paragraph{Comment 5.21}
\revcomment{The main acquisition script does not reproduce the contact subset.
\texttt{runAllDataCollection.bat} starts FBG, ATI, OptiTrack, four Mark-10 processes, motors,
and an undocumented gyroscope, but does not start the Resense/HEX12 acquisition process required
for the contact experiments.}

\paragraph{Response to Comment 5.21}
\todoresp

\paragraph{Changes made to the article}
\todoresp

\paragraph{Comment 5.22}
\revcomment{The released code is not yet a fully reproducible standalone processing pipeline.
Important parameters and paths are hard-coded in \texttt{process\_data.m}, no complete top-level
README or dependency file is provided, and several required software/hardware dependencies are
not documented. A clean, parameterized Raw $\rightarrow$ Processed workflow should be provided.
The manuscript and the released processing code describe different OptiTrack calibration
procedures. Section~2.6.3 states that the per-disk marker offsets are determined once from
\texttt{references/straight\_config}, whereas \texttt{process\_data.m} (lines 141--195)
recalculates these offsets independently from the first \SI{3}{\second} of every recording. The
authors should clarify which procedure was actually used to generate the released processed
data.}

\paragraph{Response to Comment 5.22}
\todoresp

\paragraph{Changes made to the article}
\todoresp

\paragraph{Comment 5.23}
\revcomment{The description of the circle and Lissajous trajectories does not fully match the
acquisition code. The manuscript describes continuous sinusoidal commands, whereas
\texttt{read4MotorCircle.py} implements the circle using discrete target configurations and the
Lissajous trajectory using 16 sampled waypoints that are sequentially executed by position
control. The actual actuation protocol should be described accurately.}

\paragraph{Response to Comment 5.23}
\todoresp

\paragraph{Changes made to the article}
\todoresp

\paragraph{Comment 5.24}
\revcomment{The quasi-static acquisition protocol cannot be reproduced from the released
acquisition code. The manuscript describes small incremental sweeps through near-equilibrium
configurations, while the available code mainly contains two extreme target configurations and
no explicit equilibrium criterion. Either the actual acquisition script should be released or
the Methods section should be revised accordingly.}

\paragraph{Response to Comment 5.24}
\todoresp

\paragraph{Changes made to the article}
\todoresp

\paragraph{Comment 5.25}
\revcomment{The temporal synchronization discussion appears too strong. Although all sensor
processes use the same host clock, their timestamps correspond to different events such as
serial-response completion, DAQ block reads, sequential CAN queries, network reception, or
processed FBG-frame reception. A common OS clock therefore does not necessarily imply a common
physical measurement time.}

\paragraph{Response to Comment 5.25}
\todoresp

\paragraph{Changes made to the article}
\todoresp

\paragraph{Comment 5.26}
\revcomment{The manuscript states that cross-correlation was omitted for tendon tension and base
force/torque signals, while the released synchronization code nevertheless computes
motor--tension and motor--ATI correlations. The authors should clarify whether these were
exploratory analyses only and align the released validation code with the description in the
paper.}

\paragraph{Response to Comment 5.26}
\todoresp

\paragraph{Changes made to the article}
\todoresp

\paragraph{Comment 5.27}
\revcomment{The FBG--OptiTrack comparison should be described more carefully as a cross-modal
consistency check after frame registration. Since OptiTrack data are used to align the FBG frame
before the two modalities are compared, this is not a completely independent absolute
validation.}

\paragraph{Response to Comment 5.27}
\todoresp

\paragraph{Changes made to the article}
\todoresp

% =====================================================================
\subsection{Minor comments (dataset)}
% =====================================================================

\paragraph{Comment 5.28}
\revcomment{The OptiTrack validity flags (\texttt{disk\_k\_is\_valid}) are loaded in
\texttt{data\_optitrack.m} but apparently not used during processing, and the validity
information is not retained in the processed data. The handling of occluded/invalid frames
should be documented and preferably retained as a validity mask.}

\paragraph{Response to Comment 5.28}
\todoresp

\paragraph{Changes made to the article}
\todoresp

\paragraph{Comment 5.29}
\revcomment{The FBG delay value is missing from the Methods text (``a fixed offset of is
subtracted''), while the code uses approximately \SI{13.5}{\milli\second}. The numerical value
and sign convention should be stated explicitly.}

\paragraph{Response to Comment 5.29}
\todoresp

\paragraph{Changes made to the article}
\todoresp

\paragraph{Comment 5.30}
\revcomment{Delay estimation and delay correction are mixed in \texttt{process\_data.m}: the FBG
timestamps are corrected before the synchronization check is executed. The lag estimation should
preferably be performed on the uncorrected signals and separated from the subsequent correction
step.}

\paragraph{Response to Comment 5.30}
\todoresp

\paragraph{Changes made to the article}
\todoresp

\paragraph{Comment 5.31}
\revcomment{The manuscript refers to \texttt{std::chrono::high\_resolution\_clock} for FBG
timestamps, whereas the released C++ code stores timestamps using
\texttt{std::chrono::system\_clock}.}

\paragraph{Response to Comment 5.31}
\todoresp

\paragraph{Changes made to the article}
\todoresp

\paragraph{Comment 5.32}
\revcomment{All four sequentially queried motor states are stored using a single timestamp.
Either individual timestamps should be used or the total four-motor readout latency should be
quantified.}

\paragraph{Response to Comment 5.32}
\todoresp

\paragraph{Changes made to the article}
\todoresp

\paragraph{Comment 5.33}
\revcomment{The ATI timestamp is assigned after reading and averaging a five-sample DAQ block.
The timestamp therefore represents approximately the end rather than the center of the
measurement block.}

\paragraph{Response to Comment 5.33}
\todoresp

\paragraph{Changes made to the article}
\todoresp

\paragraph{Comment 5.34}
\revcomment{The manuscript specifies a \SI{2}{\second} pre-motion period, while
\texttt{read4MotorCircle.py} uses \SI{3}{\second}. This should be made consistent.}

\paragraph{Response to Comment 5.34}
\todoresp

\paragraph{Changes made to the article}
\todoresp

\paragraph{Comment 5.35}
\revcomment{The common \SI{100}{\hertz} resampling interval is based primarily on the motor
recording duration rather than the common overlap of all sensor streams. The authors should
verify that this does not introduce undefined or NaN edge samples.}

\paragraph{Response to Comment 5.35}
\todoresp

\paragraph{Changes made to the article}
\todoresp

\paragraph{Comment 5.36}
\revcomment{Processed signals are generated using zero-phase \texttt{filtfilt}, which is acausal.
This is appropriate for offline analysis but should be explicitly stated in the Usage Notes
because the dataset is also proposed for online/state-estimation applications.}

\paragraph{Response to Comment 5.36}
\todoresp

\paragraph{Changes made to the article}
\todoresp

\paragraph{Comment 5.37}
\revcomment{The reported Mark-10 noise values should be checked.
\texttt{analyze\_reference\_noise.m} appears to calculate the standard deviation from the raw
pulley reaction force, while normal processing divides these values by two to obtain tendon
tension.}

\paragraph{Response to Comment 5.37}
\todoresp

\paragraph{Changes made to the article}
\todoresp

\paragraph{Comment 5.38}
\revcomment{The ``Euclidean distance'' used for contact-position validation is calculated only in
the $xy$ plane in \texttt{compare\_ft\_sensors.m}. It should either be referred to as planar
distance or calculated in full 3-D.}

\paragraph{Response to Comment 5.38}
\todoresp

\paragraph{Changes made to the article}
\todoresp

\paragraph{Comment 5.39}
\revcomment{The FBG files referred to as ``raw'' already contain proprietary Shape CORE
reconstructed curvature, bending angle, and 3-D shape rather than raw optical FBG measurements.
The terminology should be clarified.}

\paragraph{Response to Comment 5.39}
\todoresp

\paragraph{Changes made to the article}
\todoresp

\paragraph{Comment 5.40}
\revcomment{The statement that processed FBG quantities are retained ``without further
modification'' is inaccurate because units, coordinate frames, timestamps, filtering, and
sampling are modified during post-processing.}

\paragraph{Response to Comment 5.40}
\todoresp

\paragraph{Changes made to the article}
\todoresp

\paragraph{Comment 5.41}
\revcomment{Several additional code/documentation inconsistencies should be cleaned up,
including incorrect torque units in plot labels, obsolete function calls, inconsistent comments
on FBG column layout, and incomplete MATLAB/Python dependency documentation.}

\paragraph{Response to Comment 5.41}
\todoresp

\paragraph{Changes made to the article}
\todoresp

% =====================================================================
\subsection{Closing summary and recommendation}
% =====================================================================

\paragraph{Reviewer summary}
\revcomment{Overall, the paper addresses an existing research gap in the field of soft robotics.
The primary motivation and the targeted scientific contribution are clear and valid. However,
the paper and the dataset need more thoroughness and carefulness to reliably deliver the
intended contribution and allowing others to use the dataset. A more precise use of terms and
definitions is required in the manuscript. Also, several inconsistencies between the manuscript,
the released acquisition/processing code, and the documented data structure need to be resolved.
These include the interpretation of processed motion-capture data, calibration procedures,
trajectory generation, temporal synchronization, filtering, and the reproducibility of the
released processing pipeline. The dataset appears promising, but these issues should be
corrected to ensure reliable third-party reuse.}

\paragraph{Reviewer recommendation}
\revcomment{A major revision is necessary before the paper and the dataset can be published.}

\paragraph{Response to the closing summary}
\todoresp

% =====================================================================
%  Bibliography
%  These five entries are the datasets/works Reviewer 5 asks to be cited
%  in Comment 5.1. Merge them into the bibliography of main.tex.
% =====================================================================
\begin{thebibliography}{10}

\bibitem{Hansen2026}
H.~M. Hansen, B.~K. S{\ae}b{\o}, K.~Y. Pettersen, J.~T. Gravdahl, and M.~{Di Castro},
``Data-driven dynamic modeling of a tendon-actuated continuum robot,''
{\em arXiv:2605.18720}, 2026.
DOI: \url{https://doi.org/10.48550/arXiv.2605.18720}.

\bibitem{Muhmann2025}
C.~Muhmann, R.~M. Grassmann, M.~Bartholdt, and J.~Burgner-Kahrs,
``Toward dynamic control of tendon-driven continuum robots using Clarke transform,''
{\em IEEE/RSJ International Conference on Intelligent Robots and Systems (IROS)}, 2025.
DOI: 10.1109/IROS60139.2025.11247518.

\bibitem{Zhang2025}
G.~Zhang, Z.~Chen, and L.~Wang,
``Virtual-work based shape-force sensing for continuum instruments with tension-feedback
actuation,'' {\em arXiv:2501.05418}, 2025.

\bibitem{Liu2026}
Y.~Liu, K.~Luo, Q.~Tian, and H.~Hu,
``An efficient data-driven framework of hybrid dynamics for real-time modeling and control of
continuum robots,'' {\em Mechanism and Machine Theory}, vol.~220, p.~106361, 2026.
DOI: 10.1016/j.mechmachtheory.2026.106361.

\bibitem{Wang2025}
E.~Wang, Z.~Deng, C.~Pan, B.~He, and J.~Zhang,
``Estimating continuum robot shape under external loading using spatiotemporal neural
networks,'' {\em IEEE/RSJ International Conference on Intelligent Robots and Systems (IROS)},
2025.
DOI: 10.1109/IROS60139.2025.11245975.

\end{thebibliography}

\end{document}
